% !TEX ROOT = ../ersti.tex

\vspace*{-2mm}
\section[Studierendenvertretung in der Fachschaft]{Studierendenvertretung}
\vspace*{-1mm}
Wer vertritt meine Interessen gegenüber der Universität? Verkürzt heißt die Antwort: Vertritt dich selbst!

Das Wort \emph{Studierendenvertretung} als Beschreibung für die Arbeit der Fachschaft MathPhysInfo ist ein wenig missverständlich, weil man meinen könnte, dass einmal im Jahr eine Vertretung der Studierenden gewählt wird, die dann schon die Arbeit macht. Offiziell ist das auch so, weil das Landeshochschulgesetz es so vorschreibt. Tatsächlich meint \emph{Studierendenvertretung} hier jedoch die Vertretung der Studierenden durch sich selbst -- uns alle.

\vspace*{-2mm}
\subsection{Offenheit der Fachschaft}
Deshalb passiert die meiste Arbeit in den Fachschaftssitzungen und in Arbeitskreisen. In einer Sitzung sind alle Studis des jeweiligen Fachs willkommen, redeberechtigt und können an jedem Entscheidungsprozess gleichberechtigt mitwirken.

Offiziell wird die Fachschaft auf zwei Wegen legitimiert. Zum einen legitimiert die jährliche Wahl von Fachschaftsmitgliedern in den Fakultätsrat die Arbeit der Fachschaft an und in der Fakultät. Zum anderen wird die Arbeit der Fachschaft im Rahmen der Verfassten Studierendenschaft und außerhalb der Fakultät durch die jährliche Wahl der Fachschaftsräte legitimiert, die formal bei allen Entscheidungen das letzte Wort haben. Tatsächlich zieht die Fachschaft ihre Legitimation aus ihrer basisdemokratischen Offenheit. Statt nur einmal im Jahr zur Urne zu trotten, kannst du jede Woche bei der Fachschaftssitzung selbst mitentscheiden! Die Fachschaftsräte sind an die dort getroffenen Entscheidungen gebunden.



\vspace*{-2mm}
\subsection{Gremien und Mandatierung}
Wer als Vertreterin der Fachschaft in ein Gremium gewählt wird, ist nicht nur dem eigenen Gewissen verpflichtet, sondern der ganzen Fachschaft. Diese diskutiert alle Themen, die im Gremium wahrscheinlich relevant werden, findet eine Fachschaftsposition und beauftragt die Gewählten, diese Position zu vertreten. Die Gremienvertreterinnen werden somit von der Fachschaftssitzung meistens imperativ mandatiert.

% \begin{figure*}[b]
%	\fboxsep5mm
%	\noindent\framebox{\parbox{.935\textwidth}{\input{mathphys/konsensstufen}}}
%\end{figure*}


\vspace*{-2mm}
\subsection{Arbeitskreise und Arbeitsgruppen}
In zwei großen Fakultäten mit Studis vom ersten Semester bis zur Doktorarbeit gibt es natürlich äußerst verschiedene Interessen, und jede setzt ihre eigenen Prioritäten. Das ist super, weil daraus total verschiedene Aktivitäten entstehen -- vom Spieleabend bis zur Vortragsreihe über IT-Sicherheit finden diverse Ideen im Fachschaftsumfeld engagierte Mitstreiterinnen für ihre Umsetzung.

Damit nicht die ganze Fachschaftssitzung darüber diskutieren muss, ob man neben Brettspielen auch Rollenspiele anbietet, gibt es verschiedene Arbeitskreise (AKs) der Fachschaft, die bei Bedarf in den Sitzungen berichten. Wenn eine Entscheidung oder Geld benötigt wird, entscheidet die Sitzung. Dort werden auch AK-Gründungen bekanntgegeben.

Natürlich können auch Erstis in Arbeitskreisen mitmachen, welche es so gibt findest du auf \autopageref{AKs}




%\vspace*{-2mm}
%\subsection{AK-Queer}
%In der Hochschulgruppe PR$\sqrt{-1}$DE treffen sich queere mathphysinfo Studis zum quatschen, vernetzen und gemeinsamen Feierabend feiern.

%Wir freuen uns über neue Gesichter im neuen Semester, komm doch einfach mal vorbei! Absprachen und Ankündigungen finden sich in der über den QR Code zugänglichen WhatsApp-Gruppe.

%Bei Fragen oder Ähnlichem kannst Du dich sonst auch an ak-queer@mathphys.info wenden.
%\begin{figure}[h]
    %\includegraphics*[width=\linewidth]{ak_queer_qr_fixed.png}
%\end{figure}

\vspace*{-2mm}
\subsection{Entscheidungsfindung in der Fachschaft}

%\sidebar{\parbox{3.5cm}{\input{mathphys/konsensstufen}}} %das passt leider nicht

Die verschiedenen Interessen und Prioritäten der Studis sind aber auch eine Herausforderung. Die Fachschaft lebt davon, dass die verschiedensten Studis ihre Gedanken und ihre Position einbringen und vertreten. Es gibt also a priori keine „Fachschaftsmeinung“, es gibt nur die Fülle der Einzelmeinungen, die durch alle Anwesenden wieder verändert wird. Die verschiedensten Ansichten zu einer gemeinsamen Position und Handlungsweise zusammenzubringen, ist dann die anspruchsvolle Aufgabe der Fachschaftssitzung.

Methodisch verwenden wir dabei die Abstimmung durch einfache Mehrheit. Das heißt, dass Anträge angenommen werden, wenn mehr Stimmen dafür sind als dagegen. In der Sitzung werden dafür zuerst alle angehört, die etwas zur Diskussion beizutragen haben. Angesichts dessen werden Vorschläge gemacht, wie man in der Frage vorgehen könnte. Deren Vor- und Nachteile werden abgewogen. Dann wird über einen Vorschlag zum Vorgehen abgestimmt. 

In Fachschaftsvollversammlungen (FSVV) gelten immer die Abstimmungen aller Anwesenden, unabhängig von den Fachschaftsrätinnen. In der Fachschaftsratssitzung (FSR-Sitzung), auch Fachschaftssitzung oder einfach "Sitzung", wird das gemeinsame Meinungsbild abgefragt, bevor die Fachschaftsrätinnen ihre Bindende Meinung abgeben. In der müssen sie sich an das Meinungsbild der Sitzung halten, oder sich rechtfertigen, warum sie sich nicht daran gehalten haben.

Die Diskussionen zur Entscheidungsfindungen werden in den Sitzungen mitprotokolliert, damit wir auch nach längerer Zeit noch nachverfolgen können, warum wir uns für einen Antrag entschieden haben oder eben nicht. Diese Protokolle müssen wir auch veröffentlichen, deshalb findet ihr sie auf der Website der Fachschaft.

%\vfill



\vspace{-2mm}

\subsection{Theorie und Praxis}
Vielleicht klingt die Aufteilung auf verschiedene Sitzung und verschiedene Abstimmungen erstmal unintuitiv, vielleicht hast du auch Einwände.  Wir laden dich zunächst ein, es dir einfach mal anzuschauen! Wenn dann noch Diskussionsbedarf und Kritik an der Arbeitsweise besteht, kann und soll dies unbedingt angesprochen werden. Nichts von dem, was hier beschrieben wurde, ist in Stein gemeißelt; unsere Arbeitsweise wird ständig an unsere Erfordernisse angepasst.

Wir sind uns dessen bewusst, dass partizipative Demokratie zeitlich viel Einsatz verlangt. Solange wir die Studiengänge nicht so umgestaltet kriegen, dass allen genug Zeit dafür bleibt, ist das auch ein Legitimationsproblem. Durch Transparenz über aktuelle Themen und Entscheidungen schaffen wir Abmilderung. Doch bisher bleibt unsere Öffentlichkeitsarbeit hinter unserem Anspruch zurück. Da Fachschaftsarbeit freiwillig und ehrenamtlich ist, wird eben genau das gemacht, wofür sich jemand -- ggf. mit ein bisschen Erwartungsdruck \dots -- freiwillig Zeit nimmt. Wenn du Lust darauf hast, bist du herzlich dazu eingeladen, die Öffentlichkeitsarbeit zu verbessern! Wir freuen uns und es gibt immer Leute, von denen man hier einiges lernen kann.

Vielleicht mag es auf dich so wirken, als seien alle in den Sitzungen ein eingeschworener Haufen, die eigentlich eh die gleichen Meinungen vertreten. Aber genau das versuchen wir zu verhindern, die maßgebende Mehrheit wird von jeder Person gebildet, die ihre Meinung äußert. Egal in welche Richtung deine Meinung geht, du kannst sicher sein, dass sie zu einer Diskussion und einer neuen Meinungsbildung führt. Die Fachschaft sind alle, die sich einbringen wollen, mit all ihrer Persönlichkeit.

Dominanz mancher Leute in der Sitzung ist und bleibt ein Problem. Damit nicht einzelne die ganze Diskussion bestimmen, führen wir bei Bedarf Redelisten, wobei Leute vorgezogen werden, die erst wenig gesagt haben. Wenn du Unbehagen mit Redeverhalten oder Dis\-kussions\-weise empfindest, teile das der Sitzung mit. Generell appellieren wir an unsere Reife. Und: Nachfragen ist erlaubt und erwünscht!

Wir alle sind Lernende in fairem demokratischem Verhalten.

\subsection{Feedback willkommen!}
Da du als Ersti einen recht unverklärten Blick auf die Fachschaftssitzung hast, sollst du Gelegenheit bekommen, deine Eindrücke zu äußern. Damit wir Kritik, Anregungen und Streitfragen zur Arbeitsweise der Fachschaft in Ruhe gemeinsam besprechen können, soll nach ein paar Sitzungen die Diskussion darüber in der Fachschaftssitzung geführt werden. Den Bedarf dafür kannst du einfach in einer Sitzung anmelden, dann machen wir einen Termin dafür fest.

Haben wir Dein Interesse geweckt? Gut, dann sehen wir uns am Mittwochabend um 18:00 Uhr!

\begin{center}
    \large
    \textbf{Fachschaftssitzungen}

    \textbf{jeden Mittwoch um 18 Uhr \gls{c.t.}}
\end{center}
